\subsection{讨论}\label{sec:7.4}
\paragraph{\autoref{sec:7.1}}
Wasserstein 测地线的详细讨论见\cite[第5章]{ref315}与\cite[第7章]{ref18}. 更一般地，任意长度空间上的 Wasserstein 空间也是长度空间.

\paragraph{\autoref{sec:7.2}}
Wasserstein 空间的非负曲率在\cite[第7.3节]{ref18}中证明. 更一般地，非负曲率 Alexandrov 空间上的 Wasserstein 空间也是非负弯曲的.

\autoref{exer:7.2}中的曲线称为广义\index{geodesic@测地线（geodesic）!generalized@广义形式（generalized）|hyperpage}测地线（generalized geodesic），它在 Wasserstein 梯度流理论中起重要作用，有时也用于最优传输的其他应用.

如\autoref{sec:6.1}的讨论说明所提及，有一种适用于测地空间的综合 Ricci 曲率下界概念. 自然会问：为 Wasserstein 空间配备适当测度后，这个概念能否恢复它的非负 Alexandrov 曲率？遗憾的是，文献\cite{ref87}说明，即使是实数轴上平坦的 Wasserstein 空间，也不能由此得到有限的曲率下界.

\paragraph{\autoref{sec:7.3}}
Wasserstein 空间的切锥及其与切空间的关系，见\cite[第12.4节]{ref18}.
